% ECONOMICS_PROBLEM: {"id": "J24-590", "title": "Remote work and career development", "jel_code": "J24", "source_id": "OP-0358"}
% !TeX program = xelatex
\documentclass[11pt,a4paper]{article}
\usepackage[margin=23mm]{geometry}
\usepackage{fontspec}
\setmainfont{Latin Modern Roman}
\usepackage{amsmath,amssymb,mathtools}
\usepackage{xcolor,enumitem,booktabs,longtable,array,xurl,hyperref}
\definecolor{muted}{HTML}{53636C}
\hypersetup{colorlinks=true,urlcolor=blue,linkcolor=blue}
\setlength{\parindent}{0pt}
\setlength{\parskip}{5pt}
\newcommand{\R}{\mathbb R}
\newcommand{\E}{\mathbb E}
\newcommand{\N}{\mathbb N}
\newcommand{\ind}{\mathbf 1}
\newcommand{\Var}{\operatorname{Var}}
\newcommand{\Cov}{\operatorname{Cov}}
\newcommand{\supp}{\operatorname{supp}}
\DeclareMathOperator*{\argmin}{arg\,min}
\DeclareMathOperator*{\argmax}{arg\,max}
\newcommand{\entrymeta}[3]{{\sffamily\small\color{muted}\textbf{\texttt{#1}}\quad|\quad #2\par #3\par}}
\newcommand{\Problemref}[1]{\texttt{#1}}
\newcommand{\JELref}[2]{\texttt{#2} (JEL \texttt{#1})}
\begin{document}
\section*{Remote work and career development}
\textbf{Problem ID:} \texttt{J24-590}\quad \textbf{JEL:} \texttt{J24}\par
% BEGIN ECONOMICS STATEMENT
\label{problem:OP-0358}
{\sffamily\small\color{muted}Original subject group: Labor markets and work\par}
\label{definitions:problem:30007640}
\textbf{Audit classification:} Background; proposed extension. Metadata and full paper checked. The conclusion reports Turkish call-center results and onboarding tradeoffs; long-run career mentoring and promotion questions are editorial.
\subsection*{Definitions and precise research target}
\textbf{Complete editorial problem.} Consider workers hired into specified occupations at baseline and a feasible work-location policy \(r\in\{0,h,1\}\), denoting onsite, a declared hybrid schedule, and fully remote work. Let \(Y_{it}(r)\) be earnings, \(S_{it}(r)\) independently assessed skills, and \(P_{it}(r)\) promotion status at date \(t\). Let \(X_i\) include baseline experience, occupation, and commuting constraints. Define \(\tau_Y(r,x)=E[\sum_{t=0}^{T}\delta^t\{Y_{it}(r)-Y_{it}(0)\}\mid X_i=x]\), where \(\delta\in(0,1]\) is a specified discount factor and \(T\) the follow-up horizon. Estimate analogous skill and promotion effects and determine whether mentoring, contact with experienced coworkers, or evaluation practices mediate them. Initial onboarding and subsequent work location should be varied separately. Randomized location policies or credible quasi-experiments must address worker selection, endogenous exits, and effects on coworkers. Promotion differences alone cannot distinguish learning from managerial visibility. Evaluate benefits and losses for entrants separately from experienced workers. The cited call-center findings resolve some local productivity comparisons; extrapolating those results to long-term career development and other occupations is the proposed remaining empirical task.

\textbf{Definitions and admissible scope.} Write hybrid schedules as an explicit vector of onsite weekdays rather than a fraction with many possible meanings. Let \(r\) denote assignment, not workers' chosen attendance, and \(o\) an independent onboarding assignment. Baseline workers remain in the analysis after exit: earnings include all subsequent employers, promotion is separately coded as promotion at the initial firm or any firm, and skills are assessed independently of manager visibility. Define coworker exposure \(s_i\) as a prespecified function of coworkers' assignments when interference matters. Effects then compare \(Y_{it}(r,o,s_i)\) under stated joint assignment laws. Finite-horizon discounted earnings and skill outcomes have finite means. A mediation contrast requires a separately defined mentor intervention or defended sequential assumptions; the total work-location effect alone does not identify learning. Take \(T<\infty\) to be a fixed integer, and fix the baseline cohort and weighting law. In the displayed one-argument \(Y(r)\), onboarding is held at a declared \(o_0\) and coworkers are assigned according to a declared joint location regime; when these vary, use \(Y(r,o,s)\) explicitly. Skills and promotion are separate observed outcomes with fixed scales and finite expectations. For mediation, define a mentor intervention \(m\) and controlled contrasts \(E[Y(r,m)-Y(r,m')]\), or state the extra assumptions required for a natural-effect contrast. Work-location randomization alone identifies no such mediation effect. Report outcomes after exit under the fixed cohort convention.
\subsection*{Variants and assumption changes}
\begin{enumerate}
\item Onboarding by later location (source-motivated): randomize initial onsite contact and later work location independently. Estimate their interaction and retention effects for entrants; the call-center paper already reports local starter comparisons.
\item Visibility versus learning (editorial): cross location assignment with blinded performance evaluation or standardized mentoring. Compare independently assessed skill gains with promotion effects over several years and define the population after worker exits.
\end{enumerate}
\subsection*{References and source audit}
\textbf{Support and access.} Metadata and full paper checked. The conclusion reports Turkish call-center results and onboarding tradeoffs; long-run career mentoring and promotion questions are editorial. \textbf{Locator:} Section 5, pp. 9--10.
\begin{enumerate}\raggedright
\item\hypertarget{definitions:source:30007640:1}{} Cevat Giray Aksoy; Nicholas Bloom; Steven J. Davis; Victoria Marino; Cem \"{O}zg\"{u}zel. 2025. \emph{Remote Work, Employee Mix, and Performance}. Section 5, Conclusion, printed pp. 9–10.
\url{https://docs.iza.org/dp17917.pdf}.
DOI: \url{https://doi.org/10.3386/w33851}.
\end{enumerate}

\subsection*{Common conventions: Source mathematical conventions}

These conventions apply to each entry unless explicitly replaced.

\textbf{Models and identification.} A primitive model \(m\) specifies all probability laws, feasible choices, information, equilibrium and measurement rules used by its target. The admissible class \(\mathcal M\) is fixed independently of outcomes. If \(P_O(m)\) is its observable probability law and \(\tau(m)\) the target, its identified set at observable law \(P\) is
\[
 \mathcal I_\tau(P)=\{\tau(m):m\in\mathcal M,\ P_O(m)=P\}.
\]
Point identification means that this set is a singleton. An empty set signals incompatibility of data and assumptions. Coordinate bounds need not describe the joint set. Bounds are sharp only if every retained target is attainable or arbitrarily approachable by compatible models. Finite bounds require explicit restrictions on unobserved quantities; unrestricted confounding, hidden wealth, extrapolation or arbitrary equilibrium selection can yield uninformative sets.

\textbf{Causal interventions.} A potential outcome \(Y_i(p)\) is the outcome under a complete policy regime \(p\), including timing, financing, information and market spillovers. The target population and units are fixed before treatment. With interference, use the complete assignment vector or a justified exposure mapping; individual no-interference notation alone cannot identify an economy-wide intervention. A difference of mean potential outcomes does not require the cross-world joint distribution. Conditioning on post-treatment migration, survival, enrollment or employment changes the estimand and needs separate justification.

\textbf{Statistical statements.} An estimator, confidence set, forecast or test specifies a sampling law, model class and loss/coverage criterion. Uniform \((1-\alpha)\)-coverage means \(\inf_{m\in\mathcal M}\Pr_m\{\tau(m)\in C_n\}\geq1-\alpha\), or an explicitly asymptotic version. The whole parameter space trivially has coverage, so informativeness requires a stated width, risk or power objective. A forecasting improvement is relative to a named benchmark, information set and evaluation population.

\textbf{Equilibrium and optimization.} An equilibrium comprises feasible plans, optimality, beliefs where needed, budget constraints and all market/resource clearing. The primitives or a stated benchmark define each correspondence \(\mathcal E(p;m)\). Nonemptiness is assumed or proved, never inferred just from compact policy sets. A selected-equilibrium target uses a declared selection rule; a robust target quantifies over all permitted equilibria. Use suprema or infima unless attainment is established. An \(\varepsilon\)-optimal policy is feasible and within \(\varepsilon\) of the finite optimum value. Report an empty feasible set separately.

\textbf{Welfare and accounting.} Normative utility, interpersonal normalization, social weights, numeraire, discounting and the use of public receipts are primitives. Behavioral utility need not coincide with the welfare criterion. Household welfare includes ownership income and rents once; adding profits again to compensating variation generally double-counts them. A monetary equivalent is defined by an explicit transfer and price convention, with monotonicity and range conditions. Average effects, mechanism contributions, transfers and welfare losses are different targets. Infinite sums require convergence and dynamic feasible plans require terminal or no-Ponzi conditions.

\textbf{Source versions.} A working-paper date, revision date and journal publication year are distinguished. A DOI for a journal version is not automatically the DOI of an earlier manuscript. Locators refer to the stated version and printed pages where specified. A metadata-only check does not verify a claim about a full-text conclusion. Every entry's source subsection narrows the provenance claim accordingly.
% END ECONOMICS STATEMENT
\end{document}
